% Capítulo 3
\chapter{Estado da Arte}
\label{chap:estado_arte}

% TODO: parágrafo de introdução. Este capítulo analisa, através de uma revisão
% sistemática da literatura, como os conceitos do Capítulo~\ref{chap:fundamentos}
% têm sido combinados em trabalhos recentes.

\section{Metodologia da Revisão}
\label{sec:estado_arte_metodologia}
% TODO: revisão sistemática seguindo Kitchenham e Charters (2007) e o fluxo PRISMA 2020 (Page et al., 2021).
% Revisor único: documentar as adaptações ao processo.

\subsection*{Questões de Investigação}
% TODO: tabela com as questões. Proposta inicial (a validar com a orientadora):
%  QI1 - Como têm sido aplicados agentes baseados em LLM na automação de processos
%        de trabalho de conhecimento e de documentação em contexto organizacional?
%  QI2 - Que mecanismos de supervisão humana e de controlo de autonomia têm sido
%        propostos para agentes baseados em LLM?
%  QI3 - Que mecanismos de rastreabilidade e de deteção de erros são usados nestes sistemas?
%  QI4 - Como são avaliados estes sistemas em termos de eficiência, qualidade e aceitação?

\subsection*{Estratégia de Pesquisa}
% TODO: bases de dados (Scopus, IEEE Xplore, ACM DL, Web of Science), string de pesquisa, data da pesquisa

\subsection*{Critérios de Inclusão e Exclusão}
% TODO: tabelas de critérios (CI1..., CE1...)

\subsection*{Seleção dos Estudos}
% TODO: diagrama PRISMA e avaliação de qualidade

\section{Revisão da Literatura}
\label{sec:estado_arte_revisao}
% TODO: uma subsecção por questão de investigação, cada uma terminando com uma síntese

\section{Lacunas Identificadas}
\label{sec:estado_arte_lacunas}
% TODO: síntese e lacunas que justificam o trabalho
